\documentclass[11pt,letterpaper]{article}
\PassOptionsToPackage{hyphens}{url}
\usepackage[margin=.83in,top=.8in,bottom=.8in,headheight=15pt]{geometry}
\usepackage[T1]{fontenc}
\usepackage{lmodern,amsmath,amssymb,mathtools,bm,booktabs,array,graphicx,tikz}
\usetikzlibrary{arrows.meta,positioning,calc,decorations.pathreplacing,patterns}
\definecolor{ink}{HTML}{17344D}\definecolor{accent}{HTML}{087E8B}\definecolor{pale}{HTML}{EDF5F5}\definecolor{gold}{HTML}{AD6B19}
\providecommand{\dd}{\,\mathrm{d}}\providecommand{\R}{\mathbb{R}}
\providecommand{\E}{\mathbb{E}}\providecommand{\Var}{\operatorname{Var}}
\providecommand{\Cov}{\operatorname{Cov}}\providecommand{\diag}{\operatorname{diag}}
\providecommand{\norm}[1]{\left\lVert#1\right\rVert}
\providecommand{\abs}[1]{\left|#1\right|}
\usepackage{microtype,hyperref,fancyhdr,enumitem,amsthm,sectsty,needspace}
\allsectionsfont{\raggedright}
\hypersetup{colorlinks=true,linkcolor=ink,urlcolor=accent,pdfauthor={Zan Ahmad},pdftitle={Neural-network gradients and spatial filters / Review}}
\setlength{\parindent}{0pt}\setlength{\parskip}{6pt}\linespread{1.07}
\setlist{itemsep=3pt,topsep=4pt,leftmargin=*}
\emergencystretch=2em\clubpenalty=10000\widowpenalty=10000
\pagestyle{fancy}\fancyhf{}\fancyhead[L]{\small Scientific Computing in Python / 2027}\fancyhead[R]{\small Lecture 5}
\fancyfoot[L]{\footnotesize Zan Ahmad}\fancyfoot[R]{\footnotesize\thepage}\renewcommand{\headrulewidth}{.3pt}
\newtheorem{proposition}{Proposition}\newtheorem{theorem}{Theorem}\newtheorem{lemma}{Lemma}
\theoremstyle{definition}\newtheorem{definition}{Definition}\newtheorem{example}{Example}
\numberwithin{equation}{section}
\begin{document}
\setlength{\parskip}{4pt}\setlist{itemsep=2pt,topsep=3pt,leftmargin=*}
\section*{Neural-network gradients and spatial filters --- review sheet}
Use this after attempting the worked examples and practice. Explain each claim in words before calculating.

\section*{Key ideas}
\begin{itemize}
\item Affine maps combine features; nonlinearities make compositions richer.
\item Loss specifies the fitting objective; gradients guide local parameter changes.
\item Batch and spatial axes have different meanings.
\item A shared image filter is a repeated local weighted sum.
\item Fixed filters and a trained CNN are different claims.
\item Choose fitting procedures before the final test evaluation.
\end{itemize}
\section*{Common errors and repairs}
\textbf{Any neural network trained on flattened images is a CNN.} The sample trained network is fully connected. Its separate fixed-filter demonstration has convolutional structure but no learned filters.

\textbf{A falling training loss establishes held-out performance.} Training loss and held-out evaluation answer different questions.

\textbf{A negative filter response is invalid.} Signed directional filters encode which way intensity changes.

\textbf{Perfect synthetic accuracy proves robust real-world recognition.} It only demonstrates success on the particular held-out generator distribution.

\section*{Retrieve without your notes}
\begin{itemize}
\item What is an epoch?
\item Why keep training mean and scale for inference?
\item What output shape results from an 8$\times$8 image and a 3$\times$3 valid stride 1 filter?
\end{itemize}
\section*{Readiness check}
\begin{itemize}
\item Trace the shapes in a batched affine--nonlinearity--output calculation.
\item Explain loss, gradient update, epoch and mini-batch in the actual runnable model.
\item Interpret signed responses of two fixed 3$\times$3 image filters and predict output shape.
\item Compare a trained network with held-out baselines without using test data for fitting.
\item State why synthetic test accuracy does not establish general image or clinical performance.
\end{itemize}
\end{document}
